\section{Metrics, sampling uncertainty, and computational cost}\label{supp:metrics}
The primary metrics compare individual predicted curves to deterministic references on the strictly future grid $L+1,\ldots,30$. MAE averages absolute differences. RMSE takes the square root of the corresponding mean squared difference within each prediction replicate, and the archived summary averages replicate RMSE values rather than their squares. In the competing-risks run, survival and each CIF are scored separately rather than combined into some unreported aggregate target.

For a stochastic method, the sample standard deviation of its five replicate MAEs describes sensitivity to the sampling seed with the checkpoint fixed. It does not include variation from retraining. Each checkpoint first contributes its mean over the five sampling seeds. The nine single-covariate training-run means, or the five competing-risks training-run means, are then averaged equally. Their sample standard deviation describes between-run variation and still contains finite-sampling error. Pooled sampling standard deviation is the square root of the average within-run variance across the corresponding nine or five training runs. The 45 single-covariate sampling measurements represent nine training runs; the 25 competing-risks measurements represent five. They are not 45 or 25 independently trained models. Matched sampling seeds may correlate numerical errors across training runs. These standard deviations are descriptive and are not confidence intervals. Neural direct models contribute one deterministic prediction per training run and have no within-run sampling standard deviation. Averaging the prediction curves first and then scoring would answer a different question and would generally shrink the numerical variation, so we instead average already-computed losses. The original seed-42 source records include descriptive paired-patient intervals calculated from patient losses averaged over the saved sampling seeds. These intervals are conditional on the original seed-42 checkpoints and sampled prediction sets; they are not intervals for the nine-run or five-run mean results. They also do not account for selecting a sampler or comparator after looking at the results.

Outcome-based Brier scores~\citep{graf1999} and concordance answer different questions from reference-curve error. A censoring-weighted Brier calculation requires a justified censoring model, positivity at scored times, and a tie convention that matches the data. Concordance measures ordering rather than absolute probability accuracy. Neither metric can be inferred from the MAE rankings, and the distinction matters most when a checkpoint is selected by concordance but evaluated by curve MAE.

Prediction latency depends on architecture, sampled path length, batching, sampler bookkeeping, and the estimator calculation. The recorded batch size of 512 is an inference setting, while the history-model training batch size is 256. Equal particle and trajectory counts do not imply equal computational work. The reused single-covariate prediction timings and the newer model timings were not necessarily measured in the same session, and the training runs span different GPU hardware and software versions. The original Dynamic-DeepHit backend differences for seeds 50, 51 and 52 pertain to the results retained in Section~S10. The optimized Dynamic-DeepHit fits used an NVIDIA RTX PRO 6000 Blackwell Server Edition GPU; the earlier DDH timing summaries describe the original adaptation. For competing risks, the supplied pooled timing summaries use only seeds 43, 45, 46 and 47 on the same recorded GPU model; seed-42 L4 timings are retained separately. All five runs contribute to accuracy summaries. DRSA (optimized) was fitted locally on an NVIDIA GeForce RTX 3060 using PyTorch 2.11.0 with CUDA 13.0; its per-fit elapsed times and hardware metadata are supplied with the optimized results. The recorded timings therefore do not constitute a controlled speed comparison.

For a Transformer of width $d$, recomputing a full prefix of length $t$ costs order $td^2+t^2d$ at fixed depth. Summing these passes through length $T$ gives order $T^2d^2+T^3d$ per walk, and caching can bring this down to order $Td^2+T^2d$. Recurrent caching similarly avoids repeated prefix work. Actual wall-clock comparisons have to be measured with the implemented kernels and hardware rather than read off asymptotic formulas.

% Approved Suggestion 7, option B: dispersion for main-table neural configurations.
\begingroup\small
\setlength{\tabcolsep}{4pt}
\begin{longtable}{lrrrr}
\caption{Variation across the nine single-covariate training runs for the neural configurations in Table~1 of the main text. Each ATLAS training run contributes its mean MAE over five sampling seeds at $M=10{,}000$; ATLAS rows use Sobol. Between-run SD is the sample standard deviation of the nine training-run means. Minimum and maximum refer to those means. Sampling SD is the square root of the average within-run sampling variance; it is undefined for direct predictions. CPH and the true-generator oracle are unchanged controls and are omitted here. These summaries are descriptive, not confidence intervals.}\label{tab:harddispersion}\\
\toprule
Model / method & Between SD & Minimum & Maximum & Sampling SD \\\midrule
\endfirsthead
\multicolumn{5}{l}{\tablename\ \thetable\ (continued)}\\
\toprule
Model / method & Between SD & Minimum & Maximum & Sampling SD \\\midrule
\endhead
\midrule\multicolumn{5}{r}{Continued on next page}\\\endfoot
\bottomrule\endlastfoot
\multicolumn{5}{l}{\textit{Landmark $L=0$}}\\*
Weibull-MLP & 0.000239 & 0.023829 & 0.024598 & --- \\*
DRSA (optimized) & 0.007890 & 0.110197 & 0.135464 & --- \\*
DynamicDeepHit-Lite & 0.002011 & 0.004167 & 0.010698 & --- \\*
SurvTransformer-Lite & 0.002998 & 0.006564 & 0.015891 & --- \\*
Dynamic-DeepHit (optimized) & 0.003181 & 0.024291 & 0.033932 & --- \\*
SurvTransformer & 0.004859 & 0.012103 & 0.025143 & --- \\*
SurvTransformer-Discrete & 0.002161 & 0.010735 & 0.017151 & --- \\*
ATLAS-GRU / Sobol & 0.002379 & 0.004393 & 0.010309 & 0.000055 \\*
ATLAS-RNN / Sobol & 0.002659 & 0.003993 & 0.011482 & 0.000052 \\*
ATLAS-Transformer / Sobol & 0.004954 & 0.003316 & 0.018088 & 0.000064 \\
\addlinespace
\multicolumn{5}{l}{\textit{Landmark $L=20$}}\\*
Weibull-MLP & 0.000449 & 0.136597 & 0.138008 & --- \\*
DRSA (optimized) & 0.004728 & 0.140973 & 0.152804 & --- \\*
DynamicDeepHit-Lite & 0.001936 & 0.139121 & 0.145121 & --- \\*
SurvTransformer-Lite & 0.002338 & 0.139655 & 0.146120 & --- \\*
Dynamic-DeepHit (optimized) & 0.008377 & 0.037261 & 0.065852 & --- \\*
SurvTransformer & 0.009822 & 0.037009 & 0.067081 & --- \\*
SurvTransformer-Discrete & 0.004971 & 0.039368 & 0.053463 & --- \\*
ATLAS-GRU / Sobol & 0.003520 & 0.005468 & 0.016233 & 0.000072 \\*
ATLAS-RNN / Sobol & 0.004691 & 0.008798 & 0.023873 & 0.000056 \\*
ATLAS-Transformer / Sobol & 0.001328 & 0.006616 & 0.011309 & 0.000058 \\
\addlinespace
\multicolumn{5}{l}{\textit{Landmark $L=23$}}\\*
Weibull-MLP & 0.000381 & 0.185343 & 0.186221 & --- \\*
DRSA (optimized) & 0.004006 & 0.202487 & 0.214769 & --- \\*
DynamicDeepHit-Lite & 0.002553 & 0.200399 & 0.208035 & --- \\*
SurvTransformer-Lite & 0.004553 & 0.196337 & 0.212297 & --- \\*
Dynamic-DeepHit (optimized) & 0.009616 & 0.046243 & 0.074294 & --- \\*
SurvTransformer & 0.013379 & 0.067916 & 0.102157 & --- \\*
SurvTransformer-Discrete & 0.006843 & 0.048906 & 0.073837 & --- \\*
ATLAS-GRU / Sobol & 0.003435 & 0.008747 & 0.018402 & 0.000050 \\*
ATLAS-RNN / Sobol & 0.003054 & 0.013198 & 0.022547 & 0.000051 \\*
ATLAS-Transformer / Sobol & 0.002452 & 0.009737 & 0.017807 & 0.000046 \\
\addlinespace
\multicolumn{5}{l}{\textit{Landmark $L=24$}}\\*
Weibull-MLP & 0.000833 & 0.149664 & 0.152170 & --- \\*
DRSA (optimized) & 0.004159 & 0.151963 & 0.163361 & --- \\*
DynamicDeepHit-Lite & 0.001884 & 0.151875 & 0.158396 & --- \\*
SurvTransformer-Lite & 0.003835 & 0.149574 & 0.163230 & --- \\*
Dynamic-DeepHit (optimized) & 0.008846 & 0.051325 & 0.075680 & --- \\*
SurvTransformer & 0.008680 & 0.059305 & 0.083381 & --- \\*
SurvTransformer-Discrete & 0.010842 & 0.056364 & 0.091665 & --- \\*
ATLAS-GRU / Sobol & 0.004328 & 0.008064 & 0.020148 & 0.000056 \\*
ATLAS-RNN / Sobol & 0.004282 & 0.006977 & 0.019457 & 0.000059 \\*
ATLAS-Transformer / Sobol & 0.002933 & 0.007748 & 0.015835 & 0.000049 \\
\addlinespace
\end{longtable}
\endgroup


Table~\ref{tab:hard24_checkpoint_scores} gives the individual training-run scores for the common Sobol comparison. At baseline, the median MAEs across training runs are 0.006930 for ATLAS-GRU, 0.006546 for ATLAS-RNN and 0.007273 for ATLAS-Transformer. The corresponding ranges are 0.004393--0.010309, 0.003993--0.011482 and 0.003316--0.018088. Their between-run standard deviations are about 43, 51 and 77 times their pooled within-run sampling standard deviations, respectively. At this sampling budget, differences between fitted models are therefore much larger than the numerical variation from repeating the sampler.

Transformer training seeds 43 and 51 have the largest baseline errors, 0.015513 and 0.018088, but there is no single uniformly poor training run. For example, seed 42 has the smallest Transformer baseline MAE (0.003316) and its largest MAE at $L=23$ (0.017807). All nine runs remain in the analysis; these descriptive comparisons are not used to discard checkpoints.

% Approved Suggestion 14, option B: individual training-run scores.
\begingroup\small
\setlength{\tabcolsep}{4pt}
\begin{longtable}{rrrrlr}
\caption{Individual training-run survival MAEs for the single-covariate experiment. The GRU, RNN and Transformer columns refer to ATLAS with Sobol; each value averages the same five sampling seeds at $M=10{,}000$. The last two columns identify the lowest-error fitted direct comparator within that training run and landmark. The comparator may therefore differ between rows. All nine training runs are retained.}\label{tab:hard24_checkpoint_scores}\\
\toprule
Seed & GRU & RNN & Transformer & Best direct model & MAE \\\midrule
\endfirsthead
\multicolumn{6}{l}{\tablename\ \thetable\ (continued)}\\
\toprule
Seed & GRU & RNN & Transformer & Best direct model & MAE \\\midrule
\endhead
\midrule\multicolumn{6}{r}{Continued on next page}\\\endfoot
\bottomrule\endlastfoot
\multicolumn{6}{l}{\textit{Landmark $L=0$}}\\*
42 & 0.010119 & 0.008393 & 0.003316 & DynamicDeepHit-Lite & 0.010698 \\*
43 & 0.009730 & 0.005007 & 0.015513 & SurvTransformer-Lite & 0.006564 \\*
44 & 0.010309 & 0.006546 & 0.006588 & DynamicDeepHit-Lite & 0.008745 \\*
45 & 0.006742 & 0.005418 & 0.008925 & DynamicDeepHit-Lite & 0.010660 \\*
50 & 0.005138 & 0.011482 & 0.005617 & DynamicDeepHit-Lite & 0.007758 \\*
51 & 0.004393 & 0.007287 & 0.018088 & SurvTransformer-Lite & 0.007031 \\*
52 & 0.004549 & 0.005528 & 0.007273 & SurvTransformer-Lite & 0.007881 \\*
62 & 0.006930 & 0.003993 & 0.009394 & DynamicDeepHit-Lite & 0.004167 \\*
63 & 0.008425 & 0.011111 & 0.004648 & DynamicDeepHit-Lite & 0.007353 \\
\addlinespace
\multicolumn{6}{l}{\textit{Landmark $L=20$}}\\*
42 & 0.005468 & 0.013813 & 0.009060 & SurvTransformer-Discrete & 0.042674 \\*
43 & 0.008250 & 0.008798 & 0.009060 & SurvTransformer & 0.041463 \\*
44 & 0.007018 & 0.010462 & 0.008369 & SurvTransformer-Discrete & 0.044347 \\*
45 & 0.010287 & 0.009300 & 0.009128 & SurvTransformer & 0.037009 \\*
50 & 0.016233 & 0.013494 & 0.010128 & SurvTransformer-Discrete & 0.039368 \\*
51 & 0.009268 & 0.013098 & 0.009674 & SurvTransformer-Discrete & 0.052113 \\*
52 & 0.010414 & 0.015662 & 0.006616 & SurvTransformer-Discrete & 0.042726 \\*
62 & 0.006785 & 0.009518 & 0.011309 & Dynamic-DeepHit (optimized) & 0.052988 \\*
63 & 0.014140 & 0.023873 & 0.008017 & SurvTransformer-Discrete & 0.044037 \\
\addlinespace
\multicolumn{6}{l}{\textit{Landmark $L=23$}}\\*
42 & 0.008776 & 0.022547 & 0.017807 & Dynamic-DeepHit (optimized) & 0.052758 \\*
43 & 0.014619 & 0.013229 & 0.010707 & SurvTransformer-Discrete & 0.057438 \\*
44 & 0.010653 & 0.015780 & 0.012856 & SurvTransformer-Discrete & 0.053347 \\*
45 & 0.014978 & 0.013939 & 0.012150 & Dynamic-DeepHit (optimized) & 0.055410 \\*
50 & 0.015540 & 0.014584 & 0.010352 & Dynamic-DeepHit (optimized) & 0.051255 \\*
51 & 0.018402 & 0.014542 & 0.009737 & SurvTransformer-Discrete & 0.056618 \\*
52 & 0.016420 & 0.016393 & 0.012481 & Dynamic-DeepHit (optimized) & 0.046243 \\*
62 & 0.008747 & 0.013198 & 0.014520 & Dynamic-DeepHit (optimized) & 0.056692 \\*
63 & 0.014685 & 0.018786 & 0.012113 & SurvTransformer-Discrete & 0.048906 \\
\addlinespace
\multicolumn{6}{l}{\textit{Landmark $L=24$}}\\*
42 & 0.008064 & 0.019457 & 0.014011 & Dynamic-DeepHit (optimized) & 0.054247 \\*
43 & 0.010115 & 0.006977 & 0.008518 & SurvTransformer-Discrete & 0.066654 \\*
44 & 0.010159 & 0.008110 & 0.015835 & SurvTransformer-Discrete & 0.056364 \\*
45 & 0.015913 & 0.007499 & 0.009553 & SurvTransformer & 0.059305 \\*
50 & 0.020148 & 0.008532 & 0.010934 & Dynamic-DeepHit (optimized) & 0.069270 \\*
51 & 0.017121 & 0.015504 & 0.007777 & SurvTransformer & 0.063273 \\*
52 & 0.012868 & 0.012163 & 0.007748 & Dynamic-DeepHit (optimized) & 0.051325 \\*
62 & 0.009980 & 0.010619 & 0.011982 & SurvTransformer-Discrete & 0.062187 \\*
63 & 0.018180 & 0.014563 & 0.008086 & Dynamic-DeepHit (optimized) & 0.056427 \\
\addlinespace
\end{longtable}
\endgroup



For competing risks, Table~\ref{tab:crdispersion} reports training-run dispersion and Table~\ref{tab:cr_checkpoint_scores} gives individual-run scores for the common Sobol comparison. At baseline, the survival MAEs are $0.007902\pm.000969$ for ATLAS-GRU and $0.010378\pm.004211$ for ATLAS-Transformer, where the standard deviations describe between-run variation. The corresponding ranges are 0.006763--0.009192 and 0.005860--0.016214; pooled within-run sampling standard deviations are 0.000062 and 0.000080. The numerical variation from repeating Sobol sampling is therefore much smaller than the variation among these fitted checkpoints.

Transformer training seed 45 has the largest baseline survival and Death B errors, whereas seed 46 has the largest Death A errors at all three landmarks and the largest survival error at landmark 10. These observations describe variation in the saved fits; they do not establish its cause or justify discarding a run. All five runs remain in the analysis.

% Approved S46: five-run CR3 dispersion for common-Sobol main-table configurations.
\begingroup\small
\setlength{\tabcolsep}{4pt}
\begin{longtable}{lrrrr}
\caption{Variation across five competing-risks training runs for the neural configurations in Tables~2 to 4 of the main text. ATLAS uses Sobol; each training run contributes its mean MAE over five sampling seeds at $M=10{,}000$. Between SD is the sample standard deviation across the five training-run means. Minimum and maximum refer to these means. Sampling SD is the square root of the average within-run sampling variance and is undefined for direct predictions. CPH and the reused true-generator oracle are omitted. These are descriptive standard deviations, not confidence intervals.}\label{tab:crdispersion}\\
\toprule
Model / method & Between SD & Minimum & Maximum & Sampling SD \\\midrule
\endfirsthead
\multicolumn{5}{l}{\tablename\ \thetable\ (continued)}\\
\toprule
Model / method & Between SD & Minimum & Maximum & Sampling SD \\\midrule
\endhead
\midrule\multicolumn{5}{r}{Continued on next page}\\\endfoot
\bottomrule\endlastfoot
\multicolumn{5}{l}{\textit{Overall survival, landmark $L=0$}}\\*
Weibull-MLP & 0.000941 & 0.014796 & 0.017149 & --- \\*
DRSA (optimized) & 0.008051 & 0.073429 & 0.095392 & --- \\*
DynamicDeepHit-Lite & 0.002235 & 0.019689 & 0.024970 & --- \\*
SurvTransformer-Lite & 0.004247 & 0.020683 & 0.029893 & --- \\*
Dynamic-DeepHit (optimized) & 0.001566 & 0.014159 & 0.018062 & --- \\*
SurvTransformer & 0.006237 & 0.015549 & 0.030736 & --- \\*
SurvTransformer-Discrete & 0.005740 & 0.016079 & 0.030699 & --- \\*
ATLAS-GRU / Sobol & 0.000969 & 0.006763 & 0.009192 & 0.000062 \\*
ATLAS-Transformer / Sobol & 0.004211 & 0.005860 & 0.016214 & 0.000080 \\
\addlinespace
\multicolumn{5}{l}{\textit{Overall survival, landmark $L=10$}}\\*
Weibull-MLP & 0.000419 & 0.089528 & 0.090527 & --- \\*
DRSA (optimized) & 0.001908 & 0.093233 & 0.098538 & --- \\*
DynamicDeepHit-Lite & 0.000970 & 0.090057 & 0.092382 & --- \\*
SurvTransformer-Lite & 0.000527 & 0.090191 & 0.091504 & --- \\*
Dynamic-DeepHit (optimized) & 0.001064 & 0.021463 & 0.024414 & --- \\*
SurvTransformer & 0.003960 & 0.034238 & 0.044783 & --- \\*
SurvTransformer-Discrete & 0.006668 & 0.032131 & 0.048810 & --- \\*
ATLAS-GRU / Sobol & 0.000603 & 0.008022 & 0.009545 & 0.000061 \\*
ATLAS-Transformer / Sobol & 0.003193 & 0.009096 & 0.016585 & 0.000088 \\
\addlinespace
\multicolumn{5}{l}{\textit{Overall survival, landmark $L=20$}}\\*
Weibull-MLP & 0.000819 & 0.084534 & 0.086462 & --- \\*
DRSA (optimized) & 0.001269 & 0.075799 & 0.079182 & --- \\*
DynamicDeepHit-Lite & 0.001231 & 0.076062 & 0.078888 & --- \\*
SurvTransformer-Lite & 0.000742 & 0.075834 & 0.077457 & --- \\*
Dynamic-DeepHit (optimized) & 0.002024 & 0.020137 & 0.025122 & --- \\*
SurvTransformer & 0.006152 & 0.043631 & 0.059791 & --- \\*
SurvTransformer-Discrete & 0.006717 & 0.040613 & 0.057940 & --- \\*
ATLAS-GRU / Sobol & 0.000962 & 0.005345 & 0.007767 & 0.000049 \\*
ATLAS-Transformer / Sobol & 0.002568 & 0.006094 & 0.012472 & 0.000047 \\
\addlinespace
\multicolumn{5}{l}{\textit{Death A cumulative incidence, landmark $L=0$}}\\*
Weibull-MLP & 0.000326 & 0.011411 & 0.012281 & --- \\*
DRSA (optimized) & 0.008298 & 0.044096 & 0.064585 & --- \\*
DynamicDeepHit-Lite & 0.001934 & 0.015571 & 0.020233 & --- \\*
SurvTransformer-Lite & 0.005961 & 0.014263 & 0.026971 & --- \\*
Dynamic-DeepHit (optimized) & 0.001721 & 0.009296 & 0.013874 & --- \\*
SurvTransformer & 0.006619 & 0.013729 & 0.028147 & --- \\*
SurvTransformer-Discrete & 0.005366 & 0.015126 & 0.027732 & --- \\*
ATLAS-GRU / Sobol & 0.000316 & 0.005458 & 0.006254 & 0.000085 \\*
ATLAS-Transformer / Sobol & 0.002202 & 0.004788 & 0.010554 & 0.000097 \\
\addlinespace
\multicolumn{5}{l}{\textit{Death A cumulative incidence, landmark $L=10$}}\\*
Weibull-MLP & 0.000295 & 0.091250 & 0.092017 & --- \\*
DRSA (optimized) & 0.001538 & 0.094405 & 0.097997 & --- \\*
DynamicDeepHit-Lite & 0.000848 & 0.091948 & 0.093940 & --- \\*
SurvTransformer-Lite & 0.001127 & 0.091304 & 0.094065 & --- \\*
Dynamic-DeepHit (optimized) & 0.002066 & 0.015511 & 0.021124 & --- \\*
SurvTransformer & 0.005751 & 0.027236 & 0.041741 & --- \\*
SurvTransformer-Discrete & 0.004956 & 0.025591 & 0.038542 & --- \\*
ATLAS-GRU / Sobol & 0.000816 & 0.006142 & 0.008229 & 0.000072 \\*
ATLAS-Transformer / Sobol & 0.002119 & 0.006648 & 0.012090 & 0.000061 \\
\addlinespace
\multicolumn{5}{l}{\textit{Death A cumulative incidence, landmark $L=20$}}\\*
Weibull-MLP & 0.000749 & 0.073952 & 0.076010 & --- \\*
DRSA (optimized) & 0.001191 & 0.067890 & 0.070812 & --- \\*
DynamicDeepHit-Lite & 0.001343 & 0.068056 & 0.071293 & --- \\*
SurvTransformer-Lite & 0.000892 & 0.068454 & 0.070783 & --- \\*
Dynamic-DeepHit (optimized) & 0.001854 & 0.014048 & 0.018435 & --- \\*
SurvTransformer & 0.006075 & 0.030907 & 0.046494 & --- \\*
SurvTransformer-Discrete & 0.004897 & 0.030069 & 0.041021 & --- \\*
ATLAS-GRU / Sobol & 0.000581 & 0.003901 & 0.005467 & 0.000045 \\*
ATLAS-Transformer / Sobol & 0.001945 & 0.004011 & 0.008758 & 0.000040 \\
\addlinespace
\multicolumn{5}{l}{\textit{Death B cumulative incidence, landmark $L=0$}}\\*
Weibull-MLP & 0.000737 & 0.005582 & 0.007418 & --- \\*
DRSA (optimized) & 0.004795 & 0.029774 & 0.041060 & --- \\*
DynamicDeepHit-Lite & 0.000678 & 0.008407 & 0.010114 & --- \\*
SurvTransformer-Lite & 0.001600 & 0.007171 & 0.011227 & --- \\*
Dynamic-DeepHit (optimized) & 0.003015 & 0.007918 & 0.014989 & --- \\*
SurvTransformer & 0.001977 & 0.009643 & 0.015115 & --- \\*
SurvTransformer-Discrete & 0.001696 & 0.010788 & 0.014838 & --- \\*
ATLAS-GRU / Sobol & 0.000898 & 0.004206 & 0.006248 & 0.000063 \\*
ATLAS-Transformer / Sobol & 0.003685 & 0.003800 & 0.012715 & 0.000074 \\
\addlinespace
\multicolumn{5}{l}{\textit{Death B cumulative incidence, landmark $L=10$}}\\*
Weibull-MLP & 0.000204 & 0.045392 & 0.045891 & --- \\*
DRSA (optimized) & 0.001435 & 0.045525 & 0.049102 & --- \\*
DynamicDeepHit-Lite & 0.000454 & 0.043682 & 0.044718 & --- \\*
SurvTransformer-Lite & 0.000566 & 0.043382 & 0.044782 & --- \\*
Dynamic-DeepHit (optimized) & 0.001836 & 0.011190 & 0.015593 & --- \\*
SurvTransformer & 0.002759 & 0.018397 & 0.024082 & --- \\*
SurvTransformer-Discrete & 0.001679 & 0.018373 & 0.022484 & --- \\*
ATLAS-GRU / Sobol & 0.000551 & 0.003839 & 0.005024 & 0.000057 \\*
ATLAS-Transformer / Sobol & 0.001359 & 0.004914 & 0.008365 & 0.000061 \\
\addlinespace
\multicolumn{5}{l}{\textit{Death B cumulative incidence, landmark $L=20$}}\\*
Weibull-MLP & 0.000153 & 0.033250 & 0.033638 & --- \\*
DRSA (optimized) & 0.000454 & 0.031955 & 0.033077 & --- \\*
DynamicDeepHit-Lite & 0.000126 & 0.032086 & 0.032402 & --- \\*
SurvTransformer-Lite & 0.000299 & 0.031806 & 0.032552 & --- \\*
Dynamic-DeepHit (optimized) & 0.001152 & 0.008779 & 0.011758 & --- \\*
SurvTransformer & 0.002111 & 0.018771 & 0.024704 & --- \\*
SurvTransformer-Discrete & 0.001577 & 0.019173 & 0.022617 & --- \\*
ATLAS-GRU / Sobol & 0.000408 & 0.002855 & 0.003898 & 0.000037 \\*
ATLAS-Transformer / Sobol & 0.000999 & 0.003938 & 0.006411 & 0.000034 \\
\addlinespace
\end{longtable}
\endgroup

% Approved S46: individual trained-checkpoint CR3 comparisons.
\begingroup\small
\setlength{\tabcolsep}{4pt}
\begin{longtable}{rrrlr}
\caption{Individual competing-risks training-run MAEs. GRU and Transformer refer to ATLAS with Sobol, averaging five sampling seeds at $M=10{,}000$. The final columns identify the lowest-error fitted direct comparator within each training run, landmark and target. All five runs are retained.}\label{tab:cr_checkpoint_scores}\\
\toprule
Seed & GRU & Transformer & Best direct model & MAE \\\midrule
\endfirsthead
\multicolumn{5}{l}{\tablename\ \thetable\ (continued)}\\
\toprule
Seed & GRU & Transformer & Best direct model & MAE \\\midrule
\endhead
\midrule\multicolumn{5}{r}{Continued on next page}\\\endfoot
\bottomrule\endlastfoot
\multicolumn{5}{l}{\textit{Overall survival, landmark $L=0$}}\\*
42 & 0.007311 & 0.005860 & Dynamic-DeepHit (optimized) & 0.014159 \\*
43 & 0.006763 & 0.010463 & Weibull-MLP & 0.017149 \\*
45 & 0.007703 & 0.016214 & Weibull-MLP & 0.015290 \\*
46 & 0.009192 & 0.012458 & Weibull-MLP & 0.015318 \\*
47 & 0.008540 & 0.006897 & Weibull-MLP & 0.014998 \\
\addlinespace
\multicolumn{5}{l}{\textit{Overall survival, landmark $L=10$}}\\*
42 & 0.008543 & 0.009096 & Dynamic-DeepHit (optimized) & 0.023364 \\*
43 & 0.008740 & 0.009979 & Dynamic-DeepHit (optimized) & 0.024414 \\*
45 & 0.008022 & 0.013370 & Dynamic-DeepHit (optimized) & 0.023185 \\*
46 & 0.009545 & 0.016585 & Dynamic-DeepHit (optimized) & 0.022880 \\*
47 & 0.009274 & 0.009605 & Dynamic-DeepHit (optimized) & 0.021463 \\
\addlinespace
\multicolumn{5}{l}{\textit{Overall survival, landmark $L=20$}}\\*
42 & 0.006173 & 0.006094 & Dynamic-DeepHit (optimized) & 0.022848 \\*
43 & 0.006890 & 0.007659 & Dynamic-DeepHit (optimized) & 0.025122 \\*
45 & 0.005345 & 0.007877 & Dynamic-DeepHit (optimized) & 0.020137 \\*
46 & 0.007345 & 0.012472 & Dynamic-DeepHit (optimized) & 0.024227 \\*
47 & 0.007767 & 0.006370 & Dynamic-DeepHit (optimized) & 0.021430 \\
\addlinespace
\multicolumn{5}{l}{\textit{Death A cumulative incidence, landmark $L=0$}}\\*
42 & 0.005870 & 0.004788 & Dynamic-DeepHit (optimized) & 0.009296 \\*
43 & 0.006254 & 0.007143 & Weibull-MLP & 0.012281 \\*
45 & 0.006104 & 0.007306 & Dynamic-DeepHit (optimized) & 0.011359 \\*
46 & 0.005458 & 0.010554 & Dynamic-DeepHit (optimized) & 0.010783 \\*
47 & 0.005703 & 0.005658 & Weibull-MLP & 0.011903 \\
\addlinespace
\multicolumn{5}{l}{\textit{Death A cumulative incidence, landmark $L=10$}}\\*
42 & 0.008229 & 0.007508 & Dynamic-DeepHit (optimized) & 0.017679 \\*
43 & 0.007381 & 0.006648 & Dynamic-DeepHit (optimized) & 0.021124 \\*
45 & 0.006142 & 0.008001 & Dynamic-DeepHit (optimized) & 0.017833 \\*
46 & 0.006468 & 0.012090 & Dynamic-DeepHit (optimized) & 0.016918 \\*
47 & 0.007033 & 0.007842 & Dynamic-DeepHit (optimized) & 0.015511 \\
\addlinespace
\multicolumn{5}{l}{\textit{Death A cumulative incidence, landmark $L=20$}}\\*
42 & 0.005104 & 0.004011 & Dynamic-DeepHit (optimized) & 0.014806 \\*
43 & 0.004904 & 0.005644 & Dynamic-DeepHit (optimized) & 0.017678 \\*
45 & 0.003901 & 0.004401 & Dynamic-DeepHit (optimized) & 0.016148 \\*
46 & 0.004777 & 0.008758 & Dynamic-DeepHit (optimized) & 0.018435 \\*
47 & 0.005467 & 0.004457 & Dynamic-DeepHit (optimized) & 0.014048 \\
\addlinespace
\multicolumn{5}{l}{\textit{Death B cumulative incidence, landmark $L=0$}}\\*
42 & 0.004206 & 0.003800 & Weibull-MLP & 0.005582 \\*
43 & 0.004357 & 0.005023 & SurvTransformer-Lite & 0.007171 \\*
45 & 0.004380 & 0.012715 & Weibull-MLP & 0.006350 \\*
46 & 0.006248 & 0.005295 & Weibull-MLP & 0.007133 \\*
47 & 0.005508 & 0.004221 & Weibull-MLP & 0.006228 \\
\addlinespace
\multicolumn{5}{l}{\textit{Death B cumulative incidence, landmark $L=10$}}\\*
42 & 0.003839 & 0.004914 & Dynamic-DeepHit (optimized) & 0.011498 \\*
43 & 0.003907 & 0.005570 & Dynamic-DeepHit (optimized) & 0.015593 \\*
45 & 0.003965 & 0.008365 & Dynamic-DeepHit (optimized) & 0.012237 \\*
46 & 0.004759 & 0.006684 & Dynamic-DeepHit (optimized) & 0.011190 \\*
47 & 0.005024 & 0.005556 & Dynamic-DeepHit (optimized) & 0.011406 \\
\addlinespace
\multicolumn{5}{l}{\textit{Death B cumulative incidence, landmark $L=20$}}\\*
42 & 0.002855 & 0.003938 & Dynamic-DeepHit (optimized) & 0.011758 \\*
43 & 0.003123 & 0.004382 & Dynamic-DeepHit (optimized) & 0.011256 \\*
45 & 0.003034 & 0.006411 & Dynamic-DeepHit (optimized) & 0.008779 \\*
46 & 0.003414 & 0.004685 & Dynamic-DeepHit (optimized) & 0.010325 \\*
47 & 0.003898 & 0.004083 & Dynamic-DeepHit (optimized) & 0.010982 \\
\addlinespace
\end{longtable}
\endgroup



\paragraph{Competing-risks sensitivity to including training seed 42.} The primary analysis retains all five full runs. A descriptive sensitivity analysis uses the four runs on the same recorded GPU model, seeds 43, 45, 46 and 47. An ATLAS configuration still has the lowest fitted mean MAE in all nine landmark--target settings; GRU leads all nine, with conditional Monte Carlo in eight and Sobol for survival at landmark 10. However, for baseline Death B, only eight of the fourteen ATLAS configurations beat the best direct comparator in the four-run means, compared with all fourteen in the five-run means. The best-configuration result is therefore retained, while the broader all-configurations comparison is more sensitive to which runs are included. This is a sensitivity check, not a reason to exclude seed 42 or a test of statistical significance. The complete comparison is supplied in the competing-risks multi-seed supplement.
